% !TEX TS-program = lualatexmk
%% %%%%%%%%%%%%%%%%%%%%%%%%%%%%%%%%%%%%%%%%%%%%%%%%%%%%%%%%%%%%%%%
%% Test-Schriften.tex -- Test der Schriftoptionen von agfa-art
%% Liegt in ./ReadMe (nur lokal); agfa-art wird ueber texmf geladen.
%% Stand: 2026-09-24
%%
%% Mit LuaLaTeX uebersetzen. Genau EINE der \usepackage-Zeilen
%% unten aktivieren (die anderen mit % auskommentieren).
%% %%%%%%%%%%%%%%%%%%%%%%%%%%%%%%%%%%%%%%%%%%%%%%%%%%%%%%%%%%%%%%%
\documentclass[ngerman,fontsize=12pt,parskip=half]{scrartcl}

%% -- agfa-art ueber den texmf-Link (~/Library/texmf/tex/latex/agfa)
%\usepackage{agfa-art}					% Times (Standard)
%\usepackage[japanese]{agfa-art}			% Times + Japanisch
%\usepackage[lucida]{agfa-art}			% Lucida
\usepackage[lucida,japanese]{agfa-art}	% Lucida + Japanisch

\addbibresource{../bib/agfa-bib.bib}

%% -- Welche Optionen sind aktiv? (fuer die Ausgabe unten)
\newcommand{\zeigeoption}[1]{\ifthenelse{\boolean{#1}}{\textbf{an}}{aus}}

\title{Test der Schriften}
\subtitle{agfa-art: Optionen \texttt{lucida} und \texttt{japanese}}
\author{AGFA}
\date{\today}

\begin{document}
\maketitle

\section{Aktive Optionen}

\begin{tabular}{ll}
	\texttt{lucida}     & \zeigeoption{lucida} \\
	\texttt{japanese}   & \zeigeoption{japanese} \\
	\texttt{libertinus} & \zeigeoption{libertinus} \\
	\texttt{lmodern}    & \zeigeoption{lmodern}
\end{tabular}

Welche Schriften tatsächlich verwendet werden, zeigt im PDF-Betrachter
\enquote{Dokumenteigenschaften -- Schriften} (Vorschau: \enquote{Informationen einblenden}).

\section{Text}

Zwölf Boxkämpfer jagen Viktor quer über den großen Sylter Deich.
\emph{Kursiv: Falsches Üben von Xylophonmusik quält jeden größeren Zwerg.}
\textbf{Fett: Franz jagt im komplett verwahrlosten Taxi quer durch Bayern.}
\textsc{Kapitälchen: Heizölrückstoßabdämpfung.}

\texttt{Schreibmaschine: \textbackslash usepackage\{agfa-art\}}

Ziffern und Zeichen: 0123456789 -- „Anführungszeichen“ \enquote{mit enquote} -- § 12, € 3,50, 5\,\% -- ß, ä, ö, ü, Ä, Ö, Ü.

\section{Mathematik}

Im Text: Für $f \in L^1(\R)$ gilt $\norm{f}_1 = \int_{\R} \abs{f(t)} \dt$, ferner $\1_A$, $\P(X)$ und $\L{E,F}$.

\[
	\sum_{n=0}^{\infty} \frac{x^n}{n!} = \eu^{x}, \qquad
	\int f \d{\mu} \leq \norm*{\frac{1}{1+t^2}}, \qquad
	\alpha, \beta, \gamma, \varepsilon, \varphi, \vartheta, \Omega
\]

\begin{theorem}[Euler]\label{thm:euler}
Es gilt $\eu^{\pi\iu} + 1 = 0$.
\end{theorem}

\section{Japanisch}

%% Nur sinnvoll mit der Option japanese; ohne sie fehlen die Zeichen.
\ifthenelse{\boolean{japanese}}{%
Das Kanzufu (官子譜) wurde von Go Seigen (呉清源) kommentiert.
Wichtige Begriffe: 碁 (Go), 定石 (J\={o}seki), 手筋 (Tesuji), 死活 (Shikatsu).

囲碁は二人で行うボードゲームです。黒と白の石を交互に置き、地の多さを競います。
}{%
\emph{Option \texttt{japanese} ist nicht aktiv -- japanischer Text wird übersprungen.}
}

\end{document}
